
% ----------------------------------------------------------------------
\section{Style examples}
% ----------------------------------------------------------------------

% ============================================================
\begin{frame}{Blocks}

  \begin{block}{Regular Block}
    This is a regular block with \alert{alerted text} and normal content.
  \end{block}

  \begin{exampleblock}{Example Block}
    This is an example block, good for encodings or definitions.
  \end{exampleblock}

  \begin{alertblock}{Alert Block}
    This is an alert block, good for warnings or key results.
  \end{alertblock}

\end{frame}

% ============================================================
\begin{frame}{Itemize \& Enumerate}

  \begin{columns}
    \column{0.5\textwidth}
      Itemize:
      \begin{itemize}
        \item First item
        \item Second item with \alert{alert}
          \begin{itemize}
            \item Nested item
            \item Another nested
          \end{itemize}
        \item Third item
      \end{itemize}

    \column{0.5\textwidth}
      Enumerate:
      \begin{enumerate}
        \item First step
        \item Second step
          \begin{enumerate}
            \item Sub-step a
            \item Sub-step b
          \end{enumerate}
        \item Third step
      \end{enumerate}
  \end{columns}

\end{frame}

% ============================================================
\begin{frame}{Columns}

  \begin{columns}[T]
    \column{0.33\textwidth}
      \begin{block}{Column 1}
        Left column content. Can hold text, lists, or code.
      \end{block}
      \begin{itemize}
        \item item one
        \item item two
      \end{itemize}

    \column{0.33\textwidth}
      \begin{exampleblock}{Column 2}
        Middle column content.
      \end{exampleblock}
      Some regular text below the block.

    \column{0.33\textwidth}
      \begin{alertblock}{Column 3}
        Right column content.
      \end{alertblock}
      \begin{enumerate}
        \item step one
        \item step two
      \end{enumerate}
  \end{columns}

\end{frame}


% ============================================================
\begin{frame}{Theorem, Definition \& Proof}

  \begin{definition}
    A \alert{stable model} of a program $P$ is a minimal model of
    the reduct $P^M$ for some interpretation $M$.
  \end{definition}

  \begin{theorem}
    For any program $P$, if $M$ is a stable model of $P$ then
    $M \models P$.
  \end{theorem}

  \begin{proof}
    By construction of the GL reduct $P^M$, every rule head
    supported by $M$ is derivable from $P^M$. Minimality
    follows from the definition of stable models. \qed
  \end{proof}

\end{frame}


\begin{frame}[fragile]{Listings Style Reference}

\tiny
\textit{default}
\begin{lstlisting}[language=clingos]
a(1). b(X) :- a(X). #show b/1.
% Comment
\end{lstlisting}


\textit{blue-dark}
\begin{lstlisting}[style=blue-dark]
a(1). b(X) :- a(X). #show b/1.
\end{lstlisting}
\textit{blue-mid}
\begin{lstlisting}[style=blue-mid]
a(1). b(X) :- a(X). #show b/1.
\end{lstlisting}
\textit{blue-light}
\begin{lstlisting}[style=blue-light]
a(1). b(X) :- a(X). #show b/1.
\end{lstlisting}
\textit{blue-outline}
\begin{lstlisting}[style=blue-outline]
a(1). b(X) :- a(X). #show b/1.
\end{lstlisting}

\end{frame}
\begin{frame}[fragile]{Listings Style Reference (cont.)}
\textit{green-dark}
\begin{lstlisting}[style=green-dark]
a(1). b(X) :- a(X). #show b/1.
\end{lstlisting}
\textit{green-light}
\begin{lstlisting}[style=green-light]
a(1). b(X) :- a(X). #show b/1.
\end{lstlisting}
\textit{green-outline}
\begin{lstlisting}[style=green-outline]
a(1). b(X) :- a(X). #show b/1.
\end{lstlisting}
\textit{console}
\begin{lstlisting}[style=console]
$ clingo example.lp
clingo version 5.5.0
Reading from example.lp
Solving...
Answer: 1
b(1)
SATISFIABLE
\end{lstlisting}


\end{frame}

\begin{frame}{Style Reference — TikZ Base Nodes}
\begin{tikzpicture}[node distance=1.5cm]

  % --- blue row
  \node[font=\tiny\itshape, text=potnavy-mid] at (1.5, 4.3) {blue-dark};
  \node[blue-dark]    at (1.5, 3.8) {node};
  \node[font=\tiny\itshape, text=potnavy-mid] at (3.2, 4.3) {blue-mid};
  \node[blue-mid]     at (3.2, 3.8) {node};
  \node[font=\tiny\itshape, text=potnavy-mid] at (4.9, 4.3) {blue-light};
  \node[blue-light]   at (4.9, 3.8) {node};
  \node[font=\tiny\itshape, text=potnavy-mid] at (6.6, 4.3) {blue-outline};
  \node[blue-outline] at (6.6, 3.8) {node};
  % --- green row
  \node[font=\tiny\itshape, text=potnavy-mid] at (1.5, 3.1) {green-dark};
  \node[green-dark]    at (1.5, 2.6) {node};
  \node[font=\tiny\itshape, text=potnavy-mid] at (3.2, 3.1) {green-light};
  \node[green-light]   at (3.2, 2.6) {node};
  \node[font=\tiny\itshape, text=potnavy-mid] at (4.9, 3.1) {green-outline};
  \node[green-outline] at (4.9, 2.6) {node};
  % --- section label styles + grouping example
  \draw[section-box] (0.3, 0.3) rectangle (9.5, 1.7);
  \node[section-label] at (0.3, 1.7) {section};
  \node[section-title]    at (1.5, 1) {Section Title};
  \node[subsection-title] at (3.5, 1) {Subsection Title};
  \node[section-info]     at (6.0, 1) {section-info text};
  % --- arrows
  \node[font=\tiny\itshape, text=potnavy-mid] at (8.5, 4.3) {arrw};
  \draw[arrw] (7.8, 3.8) -- (9.2, 3.8);
  \node[font=\tiny\itshape, text=potnavy-mid] at (8.5, 3.1) {subarrw};
  \draw[subarrw] (7.8, 2.6) -- (9.2, 2.6);
  \node[font=\tiny\itshape, text=potnavy-mid] at (8.5, 2.3) {arrw-system};
  \draw[arrw-system] (7.8, 1.8) -- (9.2, 1.8);

\end{tikzpicture}
\end{frame}
\begin{frame}{Style Reference — Tikz custom nodes}
\begin{tikzpicture}[node distance=1.5cm]
  % --- semantic nodes row 1
  \node[font=\tiny\itshape, text=potnavy-mid] at (1.5, 4.3) {encoding};
  \node[encoding]  at (1.5, 3.8) {encoding};
  \node[font=\tiny\itshape, text=potnavy-mid] at (3.2, 4.3) {encoding2};
  \node[encoding2] at (3.2, 3.8) {encoding2};
  \node[font=\tiny\itshape, text=potnavy-mid] at (4.9, 4.3) {input};
  \node[input]     at (4.9, 3.8) {input};
  \node[font=\tiny\itshape, text=potnavy-mid] at (6.6, 4.3) {output};
  \node[output]    at (6.6, 3.8) {output};
  % --- semantic nodes row 2
  \node[font=\tiny\itshape, text=potnavy-mid] at (1.5, 3.1) {system};
  \node[system]    at (1.5, 2.6) {system};
  \node[font=\tiny\itshape, text=potnavy-mid] at (3.2, 3.1) {action};
  \node[action]    at (3.2, 2.6) {action};
  % --- usage example with arrows
  \node[font=\tiny\itshape, text=potnavy-mid] at (5.5, 0.9) {usage example};
  \node[system]   (sys)  at (1.5, 0.3) {system};
  \node[encoding] (enc)  at (3.5, 0.3) {encoding};
  \node[output]   (out)  at (5.5, 0.3) {output};
  \node[action]   (act)  at (7.5, 0.3) {action};
  \draw[arrw]      (sys) -- (enc);
  \draw[arrw]      (enc) -- (out);
  \draw[arrw-system] (act) -- (out);
\end{tikzpicture}
\end{frame}
